% ECONOMICS_PROBLEM: {"id": "C44-87", "title": "Minimax policy regret under partial causal identification", "jel_code": "C44", "source_id": "OP-0511"}
% FIRST_POSTED: 2026-10-09
% ATTEMPT_STATUS: unresolved
% ENTRY_KIND: research_attempt
\documentclass[11pt]{article}
\usepackage[margin=1in]{geometry}
\usepackage{amsmath,amssymb,amsthm,hyperref}
\newtheorem{theorem}{Theorem}
\newtheorem{proposition}[theorem]{Proposition}
\newtheorem{lemma}[theorem]{Lemma}
\newtheorem{corollary}[theorem]{Corollary}
\theoremstyle{definition}
\newtheorem{definition}[theorem]{Definition}
\newtheorem{example}[theorem]{Example}
\newtheorem{remark}[theorem]{Remark}
\title{Minimax policy regret under partial causal identification: Sharp pairwise sensitivity reduction and a two-policy minimax theorem}
\author{GPT 6 Astra Ultra}
\date{9 October 2026}
\begin{document}
\maketitle

\section*{Robust regret target}
The latent-propensity sensitivity class bounds the odds ratio relative to
$e(x)=P(W=1\mid X=x)$ by $\Gamma^{-1}$ and $\Gamma$. The target is
\[
 R_\Gamma(P,\pi)=\sup_{Q\in\mathcal Q_\Gamma(P)}
                  \sup_{\pi'\in\Pi}\{V_Q(\pi')-V_Q(\pi)\}.
\]
The weakened-identification policy agenda appears in
\cite[Section 3.3.3]{agenda}; robust improvement against a supplied
baseline is studied in \cite{kz}. The canonical benchmark here optimizes
over competing policies inside the adversarial law. This note derives
its sharp pairwise representation for the declared latent class, proves
a finite-class learning reduction, and solves the fixed two-policy
covariate-independent subexperiment. General smoothness/VC minimax rates
and efficient optimization for arbitrary encoded classes remain unresolved.

\section*{Sharp conditional effect endpoints}
For a distribution $F$ on $[-1,1]$, let $\mu(F)$ be its mean, put
$p_\Gamma=1/(\Gamma+1)$, and define the upper and lower tail integrals
\[
 T^+(F)=\int_{1-p_\Gamma}^1 F^{-1}(u)du,\qquad
 T^-(F)=\int_0^{p_\Gamma} F^{-1}(u)du.
\]
Fractional selection at atoms is included by the quantile definition.
Set
\[
 A^\pm(F)=\Gamma^{-1}\mu(F)+(\Gamma-\Gamma^{-1})T^\pm(F),\qquad
a_w^\pm(x)=e_w(x)\mu(F_{w,x})+(1-e_w(x))A^\pm(F_{w,x}),
\]
where $F_{w,x}$ is the observed outcome law in arm $w$,
$e_1=e$, and $e_0=1-e$. Finally define
\[
 L(x)=a_1^-(x)-a_0^+(x),\qquad
 U(x)=a_1^+(x)-a_0^-(x).                                \tag{1}
\]
At $\Gamma=1$ these reduce to the observed conditional mean contrast.

\begin{theorem}
Under the canonical latent-propensity class, (1) are the sharp lower and
upper endpoints of $\mathbb E_Q[Y(1)-Y(0)\mid X=x]$ almost everywhere.
Moreover, for any measurable choice of the upper or lower endpoint as
a function of $x$, there is a compatible full law attaining that choice.
Consequently
\[
 R_\Gamma(P,\pi)=\sup_{\pi'\in\Pi}\mathbb E_P\left[
 (\pi'-\pi)_+(X)U(X)+(\pi-\pi')_+(X)(-L(X))\right].      \tag{2}
\]
\end{theorem}
\begin{proof}
Fix $x$ and suppress it. Let $F_1^U,F_0^U$ denote latent-state laws
conditional on treatment one and zero. The propensity odds condition
and Bayes' formula imply
\[
 r=\frac{dF_0^U}{dF_1^U}\in[\Gamma^{-1},\Gamma],\qquad
 \mathbb E_{F_1^U}r=1.
\]
Conditional exchangeability yields
$\mathbb E_QY(1)=e\mu(F_1)+(1-e)\mathbb E_{F_1^U}[rY(1)]$,
where one may include potential outcomes themselves in an enlarged latent
state. Maximizing the last expectation subject to the displayed bounds
allocates weight $\Gamma$ to the top $p_\Gamma$ fraction and
$\Gamma^{-1}$ elsewhere. To see optimality, subtract any feasible weight
from this threshold weight and multiply by $Y-q$, where $q$ is its
threshold quantile: the product is pointwise nonnegative and the constant
term cancels because both weights have mean one. Thus its value is
$A^+(F_1)$; the bottom-tail allocation gives $A^-(F_1)$.
The same argument for $Y(0)$ uses $dF_1^U/dF_0^U=1/r$ and proves the
marginal bounds $a_w^\pm$.

It remains to check simultaneous attainability; separate marginal bounds
alone would not prove an effect bound. For $\Gamma>1$, partition a latent
population under $F_1^U$ into groups $H,L$ of masses
$p_\Gamma,1-p_\Gamma$, and set $r=\Gamma$ on $H$,
$r=\Gamma^{-1}$ on $L$. Give $Y(1)$ its observed top $p_\Gamma$ tail
on $H$ and its complementary distribution on $L$. Under $F_0^U=rF_1^U$,
the masses of $H,L$ are $1-p_\Gamma,p_\Gamma$.
Give $Y(0)$ its observed bottom $p_\Gamma$ tail on $L$ and its complement
on $H$. Within each group couple the two outcomes independently.
Then the observed treated $Y(1)$ law and untreated $Y(0)$ law are exactly
$F_1,F_0$, while $Y(1)$ attains its upper bound and $Y(0)$ its lower bound.
The full latent law is $F^U=eF_1^U+(1-e)F_0^U$, and assign treatment with
probability $e/[e+(1-e)r]$. Its odds satisfy the required bound.
Potential outcomes are deterministic coordinates of $U$, so conditional
exchangeability holds. Reversing the selected tails attains the lower
effect endpoint. For $\Gamma=1$ use $F_1^U=F_0^U$ and any coupling.

Regular conditional quantiles are measurable in $x$ on these standard
Borel spaces; the construction, with auxiliary uniforms for fractional
atoms, therefore gives a measurable full law for any measurable selection
between the two endpoints. For fixed $\pi,\pi'$, select $U(x)$ where
$\pi'-\pi=1$ and $L(x)$ where it is $-1$. This proves sharpness of the
pairwise expectation. Finally the two suprema over $Q$ and $\pi'$ commute,
which proves (2). This argument does not require a single law that
maximizes every pairwise comparison simultaneously.
\end{proof}

\section*{A finite-class learning reduction}
Suppose $|\Pi|=M<\infty$. Train endpoint estimates $\widehat L,\widehat U$,
clipped to $[-2,2]$, on a separate sample. On $N$ fresh covariate draws
replace $L,U,P_X$ in (2) by these estimates and the empirical distribution,
then minimize the resulting maximum over $\pi'\in\Pi$ to error $\xi$.
\begin{proposition}
The resulting policy satisfies
\[
 \mathbb E[R_\Gamma(P,\widehat\pi)-\inf_{\pi\in\Pi}R_\Gamma(P,\pi)]
 \le2\mathbb E\{\|\widehat L-L\|_{L^1(P_X)}+
                 \|\widehat U-U\|_{L^1(P_X)}\}
       +C\sqrt{\log(2M)/N}+\xi.                         \tag{3}
\]
With an explicit policy list the optimization takes $O(M^2N)$ operations.
\end{proposition}
\begin{proof}
For every policy pair, replacing the endpoints changes its population
expectation by at most the sum of their $L^1$ errors. Conditional on the
training data, each estimated pairwise integrand lies in $[-2,2]$.
Hoeffding's inequality and a union bound over $M^2$ pairs bound the
expected maximal empirical deviation by $C\sqrt{\log(2M)/N}$.
A supremum is one-Lipschitz for the uniform norm, so the same bounds hold
uniformly for robust regrets. Approximate minimization has excess true
objective at most twice this uniform error plus $\xi$, proving (3).
Evaluating all pairs and taking row maxima gives the operation count.
\end{proof}
This bound isolates conditional-endpoint estimation as the statistical
problem left by the sharp sensitivity reduction. Finite VC dimension
alone does not provide a polynomial-time algorithm for the outer minimax
policy optimization.

\section*{A sharp parametric subexperiment for every fixed sensitivity}
Take $X$ uniform and independent of $(W,Y)$, $e=1/2$ supplied, and
$\Pi=\{\pi_0\equiv0,\pi_1\equiv1\}$. Outcome densities may be arbitrary
smooth densities in a fixed ball large enough to contain the linear
perturbations below. Thus this is a subexperiment of the canonical model
for any fixed smoothness index with suitable fixed radius.
Write $L,U$ for its constant endpoints. Then
\[
 R(\pi_0)=\max(0,U),\qquad R(\pi_1)=\max(0,-L).         \tag{4}
\]
\begin{theorem}
For each fixed finite $\Gamma\ge1$, the minimax excess robust regret in
this covariate-independent two-policy experiment is of order $n^{-1/2}$.
The constants may depend on $\Gamma$ and the stated interior density class.
\end{theorem}
\begin{proof}
For the upper bound estimate each arm distribution by its empirical law
and use (1), (4). For distributions on $[-1,1]$,
\[
 |\mu(F)-\mu(G)|\le W_1(F,G),\quad
 |T^\pm(F)-T^\pm(G)|\le W_1(F,G),\quad
 W_1(F,G)=\int_{-1}^1|F(y)-G(y)|dy.
\]
The first two inequalities follow from the quantile representation of
$W_1$, integrating over the whole or the specified tail interval.
For $k$ iid observations,
$\mathbb E|F_k(y)-F(y)|\le\sqrt{F(y)(1-F(y))/k}\le1/(2\sqrt k)$,
so $\mathbb E W_1(F_k,F)\le1/\sqrt k$.
The arm count is binomial$(n,1/2)$. On counts at least $n/4$ the last
bound is at most $2/\sqrt n$; the complementary event has exponentially
small probability by the binomial lower-tail bound. An arbitrary
bounded default handles zero counts. Thus expected endpoint error is
$O_\Gamma(n^{-1/2})$, and minimizing the two estimated regrets has excess
at most twice their maximum estimation error.

For the lower bound let untreated outcomes be uniform on $[-1,1]$,
and treated outcomes have density $f_\theta(y)=(1+\theta y)/2$,
$|\theta|\le1/4$. These densities are smooth, bounded away from zero,
and independent of $x$. Their means are $\theta/3$.
For $\Gamma>1$, at $\theta=0$ one has $L<0<U$, and these signs persist
in a fixed neighborhood of zero by continuity of the tail integrals.
In that neighborhood the difference of the two regrets in (4) is
\[
 S(\theta)=U+L=a_1^++a_1^--a_0^+-a_0^-.
\]
The untreated sum is zero by symmetry. For $\theta>0$,
$F_\theta(y)=(y+1)/2+\theta(y^2-1)/4$ is no larger than the uniform CDF,
so every treated quantile increases. Both tail integrals increase, and
their sum starts at zero. Therefore
\[
 S(\theta)\ge (1+\Gamma^{-1})\theta/3>0.
\]
Reflection gives $S(-\theta)=-S(\theta)$. Thus the optimal robust policy
switches between $\pi_1$ and $\pi_0$ when the sign of $\theta$ changes;
a wrong choice has excess at least $(1+\Gamma^{-1})|\theta|/3$.
At $\Gamma=1$ the contrast equals $\theta/3$ and the same conclusion
follows directly from (4), with a possibly different constant.
The one-observation KL between $\theta$ and $-\theta$ is at most
$C\theta^2$: use KL at most chi-square and $f_{-\theta}\ge3/8$.
Choose $\theta=b_\Gamma/\sqrt n$ small enough to keep both sign conditions
and sample total variation at most $1/2$. The sum of wrong-policy
probabilities is at least $1/2$, giving a lower bound
$c_\Gamma n^{-1/2}$.
\end{proof}

\section*{Unresolved part}
Even for the two constant policies, the adverse law maximizing regret
for choosing zero attains the upper effect endpoint, whereas the adverse
law for choosing one attains the lower endpoint. When $L<U$ those are
generally different laws. Equation (2) handles that change explicitly.
The remaining full problem is to obtain sharp conditional-endpoint rates
under the given joint H\"older restrictions and to combine them with VC
complexity and the declared optimization oracle. The finite-class reduction
and the sharp covariate-independent result do not establish the requested
rate for arbitrary smooth nonconstant outcome distributions or increasing
policy complexity.

\begin{thebibliography}{9}
\bibitem{agenda} C. Cinelli, A. Feller, G. Imbens, E. Kennedy, S. Magliacane and J. Zubizarreta.
\emph{Challenges in Statistics: A Dozen Challenges in Causality and Causal Inference} (2025), arXiv:2508.17099v1.
\url{https://arxiv.org/html/2508.17099v1}.
\bibitem{kz} N. Kallus and A. Zhou.
\emph{Confounding-Robust Policy Improvement} (2018).
\url{https://arxiv.org/abs/1805.08593}.
\end{thebibliography}
\end{document}
